% ============================================================================
% SEC04.TEX - Section 4.4: Membuat Prefab Bullet dan Asteroid
% ============================================================================

\section{Membuat Prefab}

Dalam game shooter, kita akan menembakkan banyak peluru dan menghadapi banyak asteroid. Kita tidak mungkin membuatnya satu per satu secara manual. Solusinya adalah \textbf{Prefab}.

\subsection{Membuat Bullet Prefab}

\begin{enumerate}
    \item Buat Sprite baru (Circle), beri nama \texttt{Bullet}.
    \item Ubah ukurannya menjadi kecil $(0.2, 0.2, 1)$.
    \item Beri warna kuning atau merah.
    \item Tambahkan komponen \textbf{Rigidbody 2D} (Gravity Scale = 0) dan \textbf{Circle Collider 2D} (ceklis \textit{Is Trigger}).
    \item Tarik objek \texttt{Bullet} dari Hierarchy ke folder \textbf{Prefabs} di Project Window.
    \item Hapus \texttt{Bullet} dari Hierarchy.
\end{enumerate}

\subsection{Membuat Asteroid Prefab}

\begin{enumerate}
    \item Buat Sprite baru (Hexagon atau Circle), beri nama \texttt{Asteroid}.
    \item Ubah warna menjadi abu-abu atau coklat.
    \item Tambahkan \textbf{Rigidbody 2D} (Gravity Scale = 0) dan \textbf{Circle Collider 2D}.
    \item Jadikan Prefab dengan menariknya ke folder \textbf{Prefabs}.
    \item Hapus dari Hierarchy.
\end{enumerate}

Sekarang kita memiliki cetakan (blueprint) untuk Peluru dan Asteroid yang siap dipanggil (spawn) kapan saja melalui kode.
